\documentclass[12pt]{article}
\usepackage[utf8]{inputenc}
\usepackage[T1]{fontenc}
\usepackage{hyperref}
\title{Tokelauan Language Resources for K-12 Curriculum Development}
\author{Compiled by Claude Sonnet 5 (claude-sonnet-5) for Andr\'as Kornai}
\date{}
\begin{document}
\maketitle

\section{Introduction}
Tokelauan (Gagana Tokelau) is the language of Tokelau and of a diaspora community in New Zealand. Most of the school materials found are produced by the New Zealand Ministry of Education for teaching Tokelauan as an additional language. Tokelau's own schools follow a separate curriculum, and the searches found no materials for it.

\section{Primary Educational Resources}
\begin{itemize}
\item Gagana Tokelau: The Tokelau Language Guidelines (NZ Ministry of Education, TKI Pasifika): \url{https://pasifika.tki.org.nz/Pacific-languages/Gagana-Tokelau}
\item Muakiga! An Introduction to Gagana Tokelau, an entry-level unit for Years 7 to 10 with storybooks: \url{https://pasifika.tki.org.nz/content/download/1393/8238/file/Muakiga-unit+total.pdf}
\item Announcement of a new Tokelauan language resource: \url{https://gazette.education.govt.nz/articles/new-tokelauan-language-resource-launched}
\item Tokelau Language Week page (NZ Curriculum online): \url{https://newzealandcurriculum.tahurangi.education.govt.nz/te-v-iaho-o-te-gagana-tokelau-tokelau-language-week/5637151349.p}
\item Tokelau language resource download kit (Teacher Talk): \url{https://teachertalk.org.nz/products/tokelau-language-resource-download-kit}
\end{itemize}

\section{Linguistic and Reference Materials}
\begin{itemize}
\item Matauala online dictionary (the 1986 Tokelau Dictionary and an English-to-Tokelauan word list): \url{https://www.matauala.org.nz/dictionary}
\item Wikipedia, Tokelauan language: \url{https://en.wikipedia.org/wiki/Tokelauan_language}
\end{itemize}

\section{Religious and Community Resources}
None found beyond the Matauala site above.

\section{Digital and Audio Resources}
The Muakiga! unit comes with storybooks, and the download kit is a digital resource. Details of their audio content were not checked.

\section{Curriculum Development Context}
The NZ guidelines and the Muakiga! unit support teaching Tokelauan as an additional language in Years 7 to 10. Resources for Years 11 to 13 are still in draft, and Tokelau-based curriculum materials were not found. The language is endangered, so a first-language curriculum would need further work with the Tokelau Department of Education.
\end{document}